\section{Existing Tools}
\label{sec:related:tools}

The Ho Chi Minh City bus network has a public-facing presence in six
artefacts: an official route directory, four mobile applications, and
one publicly accessible research demonstrator.
Table~\ref{tab:related:tools} summarises each artefact's category and
analytical reach.
None supports the multi-stakeholder analytical inspection the present
project proposes, although the research demonstrator does cover one
spatial-density view that the proposed catalogue also includes.

\begin{table}[!htbp]
  \centering
  \caption{Existing public-facing artefacts for the Ho Chi Minh City bus network.}
  \label{tab:related:tools}
  \footnotesize
  \renewcommand{\arraystretch}{1.3}
  \begin{tabularx}{\linewidth}{p{2.5cm} p{2.8cm} X}
    \hline
    \textbf{Artefact} & \textbf{Category} & \textbf{Analytical reach} \\
    \hline
    MCPT route directory~\cite{mcptDirectory}     & Web directory             & Lists routes, stops, and operating companies; no aggregation, ranking, or temporal view \\
    BusMap (GoMo)~\cite{busmap}                   & Mobile routing            & Live multi-leg routing with arrival estimates; adds stop reminders and a TPBank GoMo fare-card link; HS-SV student-fare verification (eligibility, not analytics) \\
    GoBus! \cite{gobus}                           & Mobile routing            & Live multi-leg routing; alternative routing model to BusMap; no aggregated network views \\
    MultiGo / MultiPass~\cite{multigo}            & Multi-modal routing       & Combines bus, metro, river bus, and bicycle in a single itinerary; bundles the MultiPass electronic-ticketing card; an upstream data source for several digital-payment fields the present project consumes through the AFC stream \\
    VinBus~\cite{vinbus}                          & Operator-specific routing & Restricted to the Vingroup electric-bus fleet; routing only, no cross-operator analytical view \\
    \texttt{bus.dev.tris.vn} demonstrator         & Research demonstrator     & Spatial-intensity heatmap on a hexagonal grid; parameter controls for cell size (200~m to 2~km), route weighting, density coefficient, and activation function; no reliability, financial, school-aware, anomaly, or per-route surface \\
    \hline
  \end{tabularx}
\end{table}

Five of the six artefacts are passenger routing tools and answer one
analytical question between them, \textit{how do I get from here to
there by bus}. The two artefacts that step beyond routing each cover a
single additional dimension: MultiGo extends routing to multi-modal
itineraries and supplies the MultiPass payment data the present project
consumes downstream; the \texttt{bus.dev.tris.vn} demonstrator covers
one specific view (the hexagonal-bin density map) that also appears in
the proposed catalogue. The demonstrator is the single closest analogue
to the proposed system's Spatial and Connectivity theme; the present
project adopts the same Gridify-style construction surveyed by
Liu~\cite{liu2022} and credits the underlying construction in
Section~\ref{sec:solution:viz:spatial}. The demonstrator covers no
other analytical view of the network.

\subsection{Gap}
\label{sec:related:tools:gap}

Taken together, the six artefacts serve the passenger as a navigator
and provide one isolated spatial-density view. The two user roles
other than the passenger (the operating company and the
public-transport authority) consequently have no analytical surface to
consult about the network. Driver-level data flows into the warehouse
of the proposed system as a data source, but no driver-facing screen
is produced, because drivers receive route assignments through
dispatch channels rather than through analytical dashboards. The
analytical capabilities the six artefacts do not provide are
enumerated below.

\begin{itemize}
  \item No aggregation of ridership across time windows such as days,
        weeks, or months; each artefact presents only the live or
        near-live state of the network.
  \item No comparison of one route against another along any
        quantitative axis; the directory lists routes side by side
        but does not rank or score them.
  \item No surfacing of operational reliability indicators derived
        from vehicle telemetry, such as on-time performance,
        schedule adherence, or trip completion.
  \item No reporting of financial efficiency per route or per
        operating company; revenue and trip-cost information is not
        visible in any of the artefacts.
  \item No visualisation of spatial demand density at the level of
        the stop or of an urban subregion; no map-based view of
        where ridership concentrates.
  \item No integration of demand data (AFC) with operational data
        (GPS) under a single analytical surface.
\end{itemize}

The present project is motivated by the absence of such a surface and
addresses the absence through the analytical platform developed in
Chapter~\ref{chapter:solution}.
